\section{Минимальные остовные деревья}

\Def Пусть $G \langle V, E \rangle$ -- граф. Тогда граф $G' \langle V', E' \rangle$, будем называть \textit{подграфом}, если $V' \subset V, E' \subset E$. Если граф $G'$ является подграфом $G$, будем обозначать это $G' \subset G$.

\Def Пусть $T\langle V_T, E_T \rangle \subset G\langle V, E \rangle$. Если $T$ -- дерево и более того $V_T = V$, тогда $T$ -- \textit{остовное дерево} графа $G$. Остовное дерево также называют остовом.

\Def Пусть $G \langle V, E \rangle$ -- граф, $w: E \to \R$ -- весовая функция. Тогда \textit{весом графа} назовём следующую величину:
\begin{equation*}
    w(G) \overset{def}{=} \sum\limits_{e \in E} w(e)
\end{equation*}

\Def Пусть $G \langle V, E \rangle$ -- связный граф, $w: E \to \R$ -- весовая функция. \textit{Минимальным остовным деревом} назовём такое остовное дерево, вес которого минимален среди всех других остовных деревьев графа $G$. Иногда минимальное остовное дерево называют \textit{миностовом}.

Собственно, задача, которую мы будем решать в текущем разделе -- это алгоритмы построения тех самых миностовов. 
Все алгоритмы, которые мы разберём, по своей сути жадные и используют в своей основе одну и ту же лемму, которую мы прямо сейчас докажем.

Далее под графом мы везде позразумеваем связный граф, если не оговорено обратное (ведь остовные деревья не в связных графах не имеют смысла).

\section{Лемма о безопасном ребре}

\Note Далее мы будем отождествлять множество рёбер $E$ с соответствующим графом $G \langle V, E \rangle$, множество вершин которого совпадает с множеством вершин, используемых в рёбрах из $E$.

Всюду далее мы считаем, что у нас задан некоторый граф $G \langle V, E \rangle$ и весовая функция $w: E \to \R$.

\Def Пусть $A$ -- некоторое множество рёбер, которое является подграфом некоторого миностова. Пусть $e \in E$. Ребро $e$ будем называть \textit{безопасным} относительно $A$, если $A \cup \{e\}$ -- подграф некоторого миностова.

Ясно, что такие ребра чрезвычайно полезны и из определения безопасного ребра уже в принципе можно сделать алгоритмы, которые строят миностов. 
Но все это упирается в то, как искать вообще такие рёбра. Следующая лемма немного нам прояснит как они ведут себя.

\Def Разрезом будем называть пару $(S, T): S \cap T = \varnothing, S \cup T = V$.

\Def Будем говорить, что ребро $(u, v) \in E$ пересекает разрез $(S, T)$, если $u \in S$ и $v \in T$, либо $u \in T$ и $v \in S$.

\Lm{О безопасном ребре} Пусть множество рёбер $A$ -- подграф некоторого миностова. Пусть $(S, T)$ -- разрез $G$ таков, что ни одно ребро из $A$ не пересекает разрез $(S, T)$.
Если $e \in E$ -- ребро минимального веса из тех, что пересекают разрез $(S, T)$, то $e$ -- безопасное ребро для $A$.
\begin{proof}
    
    Пусть $T$ -- некоторый миностов, что $A \subset T$ (такой существует, по определению $A$). Если $A \cup \{e\} \subset T$, то теорема доказана.
    Рассмотрим случай, когда $A \cup \{e\} \not\subset T$.  Будем считать, что $e = (u, v)$ и $u \in S, v \in T$.

    Заметим, что при добавлении $e$ в $T$ образуется цикл (пример этого см. на рисунке ниже). Также заметим, что  в этом цикле есть хотя бы одно ребро пересекающее разрез $(S, T)$ отличное от $(u, v)$ (на рисунке это ребро $(b, c)$).
    Обозначим данное ребро за $e'$.

    \begin{figure}
        \centering
       \begin{tikzpicture}[
            main_node/.style={circle, draw, thick, minimum size=6mm, inner sep=0pt, fill=white},
            edge_style/.style={thick},
            oval_style/.style={ellipse, draw, very thick, inner xsep=1.5cm, inner ysep=1.5cm},
            blue_line_style/.style={blue, very thick},
            highlight_edge/.style={red, line width=2mm, opacity=0.2, line cap=butt},
            label_style/.style={font=\Large\bfseries}
        ]

        \coordinate (center_left) at (-2.5, 0);
        \coordinate (center_right) at (2.5, 0);

        \node[oval_style] (oval_L) at (center_left) {};
        \node[oval_style] (oval_R) at (center_right) {};

        \node[label_style, below=0.5cm of oval_L.south] {$S$};
        \node[label_style, below=0.5cm of oval_R.south] {$T$};


        \node[main_node] (a) at ($(center_left) + (-0.5, 1)$) {$a$};
        \node[main_node] (b) at ($(center_left) + (-0.5, -0.5)$) {$b$};
        \node[main_node] (u) at ($(center_left) + (1.5, 0.5)$) {$u$};

        \node[main_node] (v) at ($(center_right) + (-1.5, 0.5)$) {$v$};
        \node[main_node] (c) at ($(center_right) + (1.0, -0.5)$) {$c$};

        \draw[edge_style] (a) -- (u);
        \draw[edge_style] (b) -- (u);
        \draw[edge_style] (b) -- (c);
        \draw[edge_style] (v) -- (c);

        \draw[blue_line_style] (u) -- (v);

        \draw[highlight_edge] ($(v)!3mm!(c)$) -- ($(c)!3mm!(v)$);
        \draw[highlight_edge] ($(c)!3mm!(b)$) -- ($(b)!3mm!(c)$);
        \draw[highlight_edge] ($(b)!3mm!(u)$) -- ($(u)!3mm!(b)$);

        \end{tikzpicture} 

        \caption{Пример цикла, который образовался при добавлении ребра $(u, v)$}
    \end{figure}

    Удалим ребро $e'$ из $T$. Как следствие $T$ разделяется на две компоненты связности. Теперь возьмем и соединим эти две компоненты ребром $e$, получив некоторое новое остовное дерево $T'$ (на примере ниже, удаляем $(b, c)$ и вставляем $(u, v)$).

    Далее, заметим, что:
    \begin{equation*}
        w(T') = w(T) - w(e') + w(e)
    \end{equation*}

    Также $e'$ -- ребро пересекающее разрез при том, что $e$ выбрано как ребро минимального веса из всех пересекающих разрез. Из этого следует, что:
    \begin{equation*}
        w(e) - w(e') \leq 0
    \end{equation*}
    Комбинируя с предыдущим равенством, получаем:
    \begin{equation*}
        w(T') = w(T) - w(e') + w(e) \leq w(T)
    \end{equation*}
    \begin{equation*}
        w(T') \leq w(T)
    \end{equation*}
    Но в силу того, что $T$ -- остовное дерево минимального веса, то $T'$ -- тоже миностов. Теперь покажем, что $A \cup \{e\}$ есть подграф $T'$.
    Заметим, что $e' \not\in A$, по определению разреза $(S, T)$. Поэтому если отождествить $T$ и $T'$ с множеством соответствующих рёбер, то получаем следующие соотношения:
    \begin{equation*}
        A \subset T \setminus \{ e' \} \Rightarrow A \cup \{ e \} \subset (T \setminus \{e'\}) \cup \{ e \} = T'
    \end{equation*}
    \begin{equation*}
        A \cup \{ e \} \subset T'
    \end{equation*}

    Получили, что $A \cup \{e\}$ есть подграф некоторого миностова $T'$, значит, по определению, $e$ -- безопасное.

\end{proof}
